% docs/math/shape_derivation_body.tex
% The frames of the shape-sensitivity derivation, shared by the presentation deck
% (shape_derivation.tex) and the documentation deck (shape_derivation_code.tex): the two
% decks match frame for frame; \codenote{...} is shown in the documentation deck only.
% Describes the estimator as implemented on main (2026-10-09).

% ---------------------------------------------------------------------------------------
\section{Problem}

\begin{frame}{Outline}
\begin{enumerate}
\item The problem: the Boltzmann equation on a body-dependent domain; a windowed drag objective
\item DSMC as a Markov chain, and every place the geometry enters it
\item The value function (adjoint), carried on the particles and pulled back through
  moves and collisions
\item The first variation of one step: pathwise, hit-reward, score, boundary and operator terms
\item The seven terms of the shape gradient, one by one
\item Assembly, the optimiser, and what is exact and what is approximate
\end{enumerate}
\vspace{1em}
Notation: step $k$, time step $\Delta t$, $F_{\mathrm{num}}$ molecules per simulator
particle, superscript $m$ = the state after the move of step $k$ (before its collisions).
\codenote{Mathematical reference note \code{docs/math/ensemble\_adjoint.tex}; validation
history and findings: \code{docs/math/shape\_validation.tex}.  Entry points:
\code{examples/ensemble\_mesh.py} (one seed's gradient), \code{examples/mesh/optimize.py},
\code{examples/mesh/optimize\_squash.py} (optimisers).  Library: \code{jax\_dsmc/shape.py}
(shape terms), \code{jax\_dsmc/ensemble\_device.py} (value backward, fits),
\code{jax\_dsmc/mesh/levelset.py} (body).}
\end{frame}

\begin{frame}{The kinetic problem}
Flow domain $\Omega(\theta) = \text{box}\setminus B(\theta)$; the body is a $C^1$ level set
\[
B = \{x : \phi(x; c) < 0\},\qquad \phi(x;c) = \sum_j c_j B_j(x),\qquad n = \nabla\phi/|\nabla\phi| ,
\]
$B_j$ cubic B-splines on knots of spacing $h_\phi$, $n$ pointing into the flow.
Boltzmann equation for the one-particle density $f(t,x,v)$ in $\Omega(\theta)$:
\[
\partial_t f + v\cdot\nabla_x f = Q(f,f),
\]
inflow (drifting Maxwellian flux) at $x = 0$, outflow at $x = L$, specular $y$ walls, and on
the body, for $v\cdot n > 0$,
\[
|v\cdot n|\, f(x, v) = \int_{v'\cdot n<0} K(v' \to v;\, n)\, |v'\cdot n|\, f(x, v')\,\dd v',
\quad K = (1-a)\,\delta(v - R_n v') + a\, q_w(v\mid n),
\]
$R_n v' = v' - 2(v'\cdot n)\,n$; $q_w$ the diffuse half-Maxwellian of the wall.
A coefficient change moves the surface along its normal by
$V_n = -\delta\phi/|\nabla\phi| = -\sum_j \delta c_j B_j/|\nabla\phi|$: every control is a
normal displacement, chained through $c(\theta)$.
\codenote{\code{mesh/levelset.py} (\code{LevelSet}, \code{value}, \code{hit\_body};
\code{levelset\_h} $= 0.0625$, 16 samples).  Controls: \code{ensemble\_mesh.levelset\_coef},
\code{PARAMETERS} (\code{radius}, \code{centre\_x}, \code{mode2}, \code{width},
\code{squash\_y}), or all coefficients (\code{optimize.py}).  Config \code{acc} $= a$.  Cases:
\code{examples/mesh/forward.py::CASES}.}
\end{frame}

\begin{frame}{The objective: windowed drag}
Mean $x$-force on the body over a window of $W$ steps:
\[
D(\theta) = \frac{1}{W\Delta t}\int_{t_1}^{t_2}\!\!\int_{\partial B}\int_{v\cdot n<0}
m\,\big(v - \vout(v)\big)_x\,|v\cdot n|\, f\,\dd v\,\dd S\,\dd t .
\]
In DSMC it is the expectation of a sum of \emph{per-hit rewards}:
\[
D = \E\Big[\sum_{k\in\text{window}}\ \sum_{\text{body hits } j \text{ of step } k} r_{k,j}\Big],
\qquad
r_{k,j} = \frac{F_{\mathrm{num}}}{W\Delta t}\, m\,\big(\vin_{k,j} - \vout_{k,j}\big)_x .
\]
For a diffuse hit $\vout$ is the drawn re-emission velocity.

\medskip
Squash study: $J(s) = D(s) - \lambda_A\,A(s)$, an ellipse with semi-axes $R$ and $sR$,
$A(s) = \pi R^2 s$ exact, $\lambda_A$ fixed so that the free-molecular optimum is $s^* = 0.35$.
Only $\dd D/\dd s$ needs the adjoint.  No state reward, no terminal reward.
\codenote{\code{ensemble\_mesh.\_drag\_reward} (per hit; window = the last \code{window}
$= 200$ of \code{steps} $= 300$ captured steps after \code{warmup} $= 700$),
\code{objective\_value} (the same sum from the surface tallies: the reported $D$).
\code{optimize\_squash.reference} (free-molecular $D$, $\dd D/\dd s$, $\dd A/\dd s$),
\code{seed\_gradient}; history rows \code{dD\_ds}, \code{J\_prime}.}
\end{frame}

% ---------------------------------------------------------------------------------------
\section{DSMC}

\begin{frame}{DSMC as a Markov chain}
State $Z_k = \{(x_i, v_i, a_i)\}_{i=1}^N$ over fixed slots ($a$ = alive).  One step:
\begin{center}
\begin{tikzpicture}[node distance=4mm, every node/.style={font=\small},
  box/.style={draw, rounded corners=2pt, inner sep=3pt, fill=gray!8}]
\node[box] (z) {$Z_k$};
\node[box, right=of z] (m) {move + walls $\mathcal W_\theta$};
\node[box, right=of m] (i) {inject};
\node[box, right=of i] (s) {sort to cells};
\node[box, right=of s] (c) {collide (NTC/VSS) $\mathcal C_\theta$};
\node[box, right=of c] (z1) {$Z_{k+1}$};
\draw[-{Stealth}] (z) -- (m); \draw[-{Stealth}] (m) -- (i); \draw[-{Stealth}] (i) -- (s);
\draw[-{Stealth}] (s) -- (c); \draw[-{Stealth}] (c) -- (z1);
\node[below=1mm of m, font=\scriptsize] {$Z_k^m$: post-move};
\end{tikzpicture}
\end{center}
\textbf{Common random numbers.} Every stage draws from a key fixed by (seed, step, stage),
so with the key fixed each stage is a deterministic function of its input: a run can be
recorded and replayed exactly, and two geometries see the same random numbers.

\medskip
\textbf{Where $\theta$ enters the forward}
\begin{enumerate}
\item[(a)] the move: hit times, hit points, normals, specular $\vout = R_n\vin$, bounce chains;
\item[(b)] the per-hit reward itself, through $\vout(n)$;
\item[(c)] the diffuse re-emission law $q_w(\cdot\mid n)$;
\item[(d)] the hit set: which particles hit, and in which step;
\item[(e)] the collisions: NTC's attempt count carries the cut-cell flow volume $V_c(\theta)$.
\end{enumerate}
\codenote{\code{simulate.\_step}; \code{simulate.run(capture=True)} $\to$ \code{EventRecord}
(post-move \code{x}, \code{v}, \code{alive}, \code{born\_face}, \code{PairRecord} per attempt
slot, \code{WallEvents}); \code{rng.key\_for}$(\kappa, k, \text{stage})$; config
\code{record="lean"}.}
\end{frame}

\begin{frame}{NTC collisions and the cut-cell volume}
In cell $c$ with $N_c$ particles and flow volume $V_c$, the attempt count is a stochastic
rounding of
\[
\bar n_{\mathrm{att},c} = \tfrac12\,N_c(N_c-1)\,\Sigma_c\,\Delta t\,F_{\mathrm{num}}/V_c ,
\]
each attempt draws a pair uniformly and accepts it with probability $q(g)/\Sigma_c$,
$q(g) = \sigma_T(g)\,g$ the VSS rate, $\Sigma_c$ a majorant; an accepted pair scatters
with the VSS law.  Hence
\[
\Pr\{\text{pair } (i,j) \text{ collides in step } k\} \approx \frac{q(g_{ij})\,\Delta t\,F_{\mathrm{num}}}{V_c}
\quad\propto\quad \frac{1}{V_c(\theta)},
\]
and a particle's expected collisions per step are
$\nu\Delta t \approx (2\bar n_{\mathrm{att},c}/N_c)\,\E[q/\Sigma_c]$.

\medskip
$V_c(\theta) = |c\cap\Omega(\theta)|$: a cut cell's collision probabilities depend on the shape.
A cell that the body splits into disconnected flow regions gets one collision cell per region.

\medskip
Kn $= \lambda_{\mathrm{mfp}}/2R = 0.5$ on the squash case: $\nu\Delta t\approx 2.5\,\%$ per step; free-molecular: $0$.
\codenote{\code{collide/ntc.py}: \code{attempt\_count} (stochastic rounding),
\code{certified\_majorant}; \code{collide/vss.py}: \code{collision\_rate}, \code{scatter}.
Cut volumes by stratified sub-sampling (\code{GridSpec.volume}); split cells
(\code{collision\_cell}, PSEUDOCODE 4.5).  Config \code{knudsen} (\code{knudsen\_density}),
\code{a\_max} (attempt slots per cell and step; 64 at Kn 0.5).}
\end{frame}

% ---------------------------------------------------------------------------------------
\section{Adjoint}

\begin{frame}{The value function}
Value of one particle at the post-move instant of step $k$:
\[
\lambda_k(x, v) = \frac{\delta D}{\delta f^m_k(x,v)}
= \E\big[\text{rewards collected from } t^m_k \text{ on} \;\big|\; \text{one particle at } (x,v)\big].
\]
For a Boltzmann gas it solves the adjoint equation, backward in time,
\[
-\partial_t\lambda - v\cdot\nabla_x\lambda = \Lstar_f\lambda + r,\qquad
(\Lstar_f\psi)(v) = \!\int\! f(w)\!\int\! q\,\big[\psi(v') + \psi(w') - \psi(v) - \psi(w)\big]\,\dd\sigma\,\dd w ,
\]
with, at a specular wall, $\lambda(x_w, \vin) = r(\vin) + \lambda(x_w, R_n\vin)$, and at a
diffuse wall the expectation over the re-emission.

\medskip
\textbf{Particle-carried form.}  The PDE is never solved.  One value per particle, at that
particle's own state, is carried backward through the recorded run:
\[
\mu_{k+1}\ \xrightarrow{\ \text{collision pull-back}\ }\ \lambda^m_k\
\xrightarrow{\ \text{move pull-back}\ }\ \mu_k ,
\]
$\mu_k$ the values at the start of step $k$, $\lambda^m_k$ at its post-move state.
\codenote{\code{ensemble\_device.backward\_virtual} $\to$ \code{Backward.lam} ($M\times N$:
$\lambda^m_k$ per captured step and slot).  The objective's value and the timing-reward
value (along-flight identity) run as one two-column sweep: \code{ensemble\_mesh.\_backwards}
(\code{columns\_out=True}).}
\end{frame}

\begin{frame}{Backward through the move}
A particle that survives the move of step $k$ keeps its value and collects the rewards of
the hits it made during the move:
\[
\mu_{k,i} = \ind\big[a^m_{k,i}\wedge\neg\,\mathrm{born}_{k,i}\big]
\Big(\lambda^m_{k,i} + \sum_j r_{k,j,i}\Big).
\]
\begin{itemize}
\item Specular walls are deterministic maps: the value moves with the particle.
\item A diffuse wall is a Markov kernel; its realised re-emission is a sample of it, so the one
  carried value estimates $\E[\lambda(\text{after})\mid\text{before}]$ without bias
  (tower property), however many bounces the step has.
\item Particles absorbed or leaving the box carry $0$; a particle born in step $k$ carries
  nothing to $t_k$ (it did not exist).
\end{itemize}
\codenote{\code{ensemble\_device.\_move\_pullback}, \code{\_step\_pullback\_virtual};
survivors \code{alive \& (born\_face < 0)}; per-hit rewards from \code{shape.replay}
(\code{HitRewards.totals()}, one total per slot and step).}
\end{frame}

\begin{frame}{Backward through the collisions: virtual partners}
The collision step applies $e^{\Delta t\,\Lstar}$ to the values; it is applied in $S$ substeps.
In each, every particle $i$ of cell $c$ draws $K$ partners $j_\kappa$ uniformly from the
other particles of $c$, and fresh VSS angles:
\[
\lambda_i \mathrel{+}= \sum_{\kappa=1}^K w_{i\kappa}\big[\Lam(v_i') + \Lam(v_{j_\kappa}') - \Lam(v_i) - \Lam(v_{j_\kappa})\big],
\qquad
w_{i\kappa} = \frac{q(g_{i\kappa})}{\Sigma_c}\,\frac{2\,n_{\mathrm{att},c}}{N_c}\,\frac{1}{KS},
\]
$\Lam$ one smooth fit of the current values in the cell (next frames).  The weight is the
expected number of collisions of $i$ in the step (attempts per particle $\times$ acceptance)
spread over $KS$ samples; all four values come from one fit, so their errors cancel.
\begin{itemize}
\item Substep $s$ fits on the velocities as they stood at that point of the NTC sweep
  (``staged'' nodes, reconstructed from the record's slot order); a particle whose
  velocity changed across a substep boundary moves its value through the fit.
\item Splitting error $O(\Delta t/S)$; noise far below one realised pair per collision.
\end{itemize}
\textbf{The collision increment} of slot $i$ in step $k$,
$\;c_{k,i} = \sum_{\text{substeps}}\sum_\kappa w_{i\kappa}\Phi_{i\kappa} \approx \Delta t\,(\Lstar\lambda)_i$,
is kept: the along-flight, order and cut-volume terms use it.
\codenote{\code{ensemble\_device.\_substep}, \code{\_partner\_pass}, \code{staged\_velocities},
\code{\_boundary\_move}; \code{ensemble\_backward.configure("virtual", S, K, "staged")},
default $(S,K) = (4,8)$, collisionless runs $(1,2)$.  \code{backward\_virtual(collision\_out=True)}
$\to$ \code{Backward.collision} ($c_{k,i}$, f32, objective column) and
\code{Backward.cell\_collision} (per-cell sums, f64).  Stability guard: rms $\lambda$ growth
$\le$ \code{GROWTH\_LIMIT} $= 100$ (\code{ensemble\_inlet\_shock.device\_backward}).}
\end{frame}

\begin{frame}{The value field of the collision pull-back: one polynomial per cell}
The pull-back needs $(3+2K)N$ field values per substep (the node, two ghosts per virtual partner, the
boundary move).  In the scaled coordinates $\tilde p = (x/\Delta x,\ v/v_{\mathrm{s}})$, $v_{\mathrm{s}} = \sqrt{kT/m}$,
each collision cell $c$ gets \textbf{one} least-squares fit of its nodes' carried values, at offsets
$o = \tilde p - \bar p_c$ from the node mean:
\[
\phi(o) = \big(1,\ o_x,\ o_v,\ \{o_{v,a}o_{v,b}\}_{a\le b},\ o_x^{2}\big),\qquad
\Big(\sum_{n\in c}\phi_n\phi_n^\top + \rho\,\bar d\,\mathrm{diag}(0,1,\dots,1)\Big)\beta_c = \sum_{n\in c}\phi_n\lambda_n ,
\]
and every query in $c$ reads $\Lam(x_q, v_q) = \phi(\tilde p_q - \bar p_c)\cdot\beta_c$.
\begin{itemize}
\item \textbf{Order guard}: only the columns of the orders whose cumulative count is $\le n_c/1.5$
  (14 columns in 2-D need 21 nodes; else a plane, else the cell mean).
\item \textbf{Next to a surface} ($\lambda$ jumps across $v\cdot n = 0$): three fits per cell -- nodes leaving,
  approaching ($v\cdot n_c \gtrless 0$, $n_c$ the cell's mean nearest-surface normal) and all; a query takes
  its side's fit, the unsplit one when no surface is near, its side is empty, or the cell's normals
  disagree ($|\sum n_i|/\#n_i < 0.9$: corners, opposed faces).
\item Exact for any field quadratic in $v$ (the per-cell fit reproduces its own monomials).
\end{itemize}
\codenote{\code{jax\_dsmc/cell\_fit.py}: \code{fit}, \code{fit\_split}, \code{cell\_normals},
\code{evaluate}; \code{ensemble\_device.\_field} / \code{\_query} (\code{\_Hybrid} for the optional per-query
fallback).  \code{ensemble\_backward.CELL\_BASIS} = \code{"cell:quad:adapt1.5:align0.9:both"} (default of
\code{configure} and every example's \code{--backward-basis}; \code{quadratic\_full} = the per-query fit).
Tests \code{tests/test\_cell\_fit.py}; record docs/PSEUDOCODE.md 9.7.2.}
\end{frame}

\begin{frame}{Why one fit per cell, and why the guard}
\begin{itemize}
\item \textbf{The collision term sees only the non-invariant part of $\lambda$}: for
  $\lambda\in\mathrm{span}(1, v, |v|^2)$ every four-term $\Lam(v_i')+\Lam(v_j')-\Lam(v_i)-\Lam(v_j)$ vanishes.  What the
  pull-back needs is the cell population's $v$-structure beyond the invariants -- a property of the
  cell, not of a neighbourhood of each query.
\item The per-query fit over the $k$ nearest nodes depended on $k$ ($k = 32\to64$: $+1.6\pm0.7\,\%$ of $\dd D/\dd s$,
  12 seeds; flat from $k = 64$, where it spans the whole cell and its weights only range over $[\frac12, 1]$)
  and cost $\propto k$.  The per-cell fit has no $k$: $C$ fits instead of $19N$.
\item A cell's values are mostly noise (the cell mean predicts held-out nodes at the noise floor).  An
  \emph{underdetermined} quadratic in a sparse cell answers the partner ghosts with extrapolated noise, and the
  backward recursion compounds it: the unguarded per-cell quadratic is unstable (one seed in 12: $+97\,\%$).
  The guard bounds the columns by the nodes.
\end{itemize}
\textbf{Measured} (Kn 0.5 squash vs CRN FD): at production density (16 seeds) guarded per-cell $-3.7\pm0.8\,\%$,
per-query $k = 32$ $-3.4\pm0.8\,\%$, a cubic-in-$v$ guarded fit $-0.7\pm0.8\,\%$; a seed $1100\to46$\,s.  Stable on the optimiser's complicated shapes, where
the per-query one-sided fit diverged.
\codenote{Evaluation of every interpolant tried (\code{mean}, \code{lin}, \code{quad}, \code{quadfull}, \code{cubicv},
\code{cubic}, \code{adapt}, \code{loo}, \code{pool}, \code{rb}, split rules, fallbacks): validation deck section V;
\code{scratch/cellfit/}.}
\end{frame}

\begin{frame}{The per-query fit: the shape gradient's slopes}
The shape terms need the slopes $\nabla_x\lambda$, $\nabla_v\lambda$ at reflected states, with node sets chosen
per query (next frame).  For a query $(x_q, v_q)$: the $k$ nearest admissible nodes in $\tilde p$, weights
$\omega_n = 1/\big(1 + (d_n/d_k)^2\big)$, the same quadratic design at $\delta = \tilde p_n - \tilde p_q$:
\[
\Big(\sum_n\omega_n\phi_n\phi_n^\top + \rho\,\bar d\,\mathrm{diag}(0,1,\dots,1)\Big)\beta
= \sum_n\omega_n\phi_n\lambda_n\qquad(\bar d\text{: mean diagonal; } \rho = 10^{-8}),
\]
\[
\Lam = \beta_0,\qquad \nabla_x\Lam = \beta_{1:1+D}/\Delta x,\qquad \nabla_v\Lam = \beta_{1+D:4+D}/v_{\mathrm{s}} .
\]
\begin{itemize}
\item Pooled over the $3^D$ block ($k = 96$): the slopes need more nodes than a value.
\item The slopes are the weak link: the value field is discontinuous next to a body (next frame).
\end{itemize}
\codenote{\code{ensemble\_device.node\_gradients} $\to$ \code{cell\_query} $\to$ \code{field\_query}
(value, $\nabla_x$, $\nabla_v$), \code{\_design} (\code{basis="quadratic\_full"}: 14 columns in 2-D),
\code{\_spd\_solve} (unrolled Cholesky), \code{RIDGE\_REL} $= \rho = 10^{-8}$; \code{cell\_nodes} (\code{npc\_max}).
Until 2026-10-11 this fit ($k = 32$) was also the collision pull-back's value field.}
\end{frame}

\begin{frame}{Which nodes: functions of the current state only}
Next to a body the post-move population splits: particles about to hit carry their
imminent hit's reward, particles that just left do not.  $\lambda$ jumps across that
boundary, and a polynomial fitted across it has a wrong slope.  The shape-gradient fits
therefore use, for a query whose last hit had normal $n$ at $x_q$:
\begin{enumerate}
\item \textbf{one-sided}: only nodes leaving that surface, $v\cdot n > 0$;
\item \textbf{pooled}: the $k = 96$ nearest such nodes from the $3^D$ block of cells around
  the query's cell (variance);
\item \textbf{half-space}: only nodes with $(x_{\mathrm{node}} - x_q)\cdot n \ge -\tfrac14\Delta x$
  (a body thinner than the block would otherwise put its far side in the pool).
\end{enumerate}
\textbf{Rule: stratify by the current state, never by the future.}  With
$\lambda = P\lambda_1 + (1-P)\lambda_0$ ($P$ the probability of some future event),
\[
\nabla\lambda = P\nabla\lambda_1 + (1-P)\nabla\lambda_0 + \nabla P\,(\lambda_1-\lambda_0);
\]
a fit inside classes defined by the future event estimates the first two terms only.
Classes defined by the present state ($v\cdot n$, position) keep all three.
\codenote{\code{ensemble\_device.node\_gradients(direction=, pool=, halfspace=)},
\code{pooled\_tables}, \code{\_halfspace\_length}.  \code{shape.VARIANTS}: the default
\code{along\_pool\_plane} = (pool 1, $k = 96$, along-flight, half-space $0.25$ cells), in
\code{ensemble\_gradient}, \code{optimize.py}, \code{optimize\_squash.py}.  Config
\code{value\_class="none"} (the future-class option \code{"hit\_ahead"} is kept for
diagnostics only and warns).}
\end{frame}

% ---------------------------------------------------------------------------------------
\section{First variation}

\begin{frame}{The first variation of one step}
Write one step as $Z' = G_\theta(Z,\omega)$, randomness $\omega\sim p_\theta$, $R_\theta(Z,\omega)$
the rewards paid at the wall during the move, $V(Z')$ the value of what follows.  Then
\begin{align*}
\frac{\dd}{\dd\theta}\E\big[R + V(Z')\big]
&= \underbrace{\E\big[\nabla V\cdot\partial_\theta G\big]}_{\text{pathwise (frozen }\omega)}
+ \underbrace{\E\big[\partial_\theta R_\theta\big]}_{\text{hit reward (frozen }\omega)}
+ \underbrace{\E\big[(R + V)\,\partial_\theta\log p_\theta(\omega)\big]}_{\text{score}}\\[-2pt]
&\quad + \underbrace{\int_{\Sigma_\theta}\jump{R + V}\, u_\Sigma\,\dd\sigma}_{\text{boundary}}
\;+\; \underbrace{\E\big[\partial_\theta(\text{collision law})\text{ at fixed particles}\big]}_{\text{operator}} ,
\end{align*}
$\Sigma_\theta$: where $G_\theta$ jumps (hit $\to$ miss; hit into the next step), $\jump{\cdot}$ the jump,
$u_\Sigma$ its speed in $\theta$.  Summed over steps with $V = \lambda$, this gives the seven terms:

\smallskip
\begin{tabular}{@{}ll@{}}
pathwise: & 1 transport of the perturbed state (with the along-flight identity)\\
hit reward, score: & 2 hit reward (tally);\quad 3 diffuse normal\\
boundary: & 4 hit $\leftrightarrow$ miss;\quad 5 window edges;\quad 6 hit/collision order\\
operator: & 7 cut-cell volume\\
\end{tabular}
\codenote{\code{shape.gradient\_variants} $\to$ \code{\_gradient} (a \code{lax.scan} over the
captured steps) $\to$ \code{\_gradient\_step}: one \code{jax.vjp} of the replayed wall map per
step, every cotangent pulled back in one batched transpose (\code{\_pullbacks}); output
\code{ShapeTerms}.  The value $\lambda$ is a cotangent, never differentiated.}
\end{frame}

\begin{frame}{The replayed wall map}
Under the fixed keys the move of step $k$ is a pure function of the geometry:
\[
(x^\pth_k,\ v^\pth_k,\ T_k,\ \mathcal E_k) = \mathcal W_\theta(x_k, v_k;\ \kappa_k),
\]
$T_k$ surface tallies, $\mathcal E_k$ the traced events ($x^h$, $\vin$, $\vout$, $n$, $t^h$ per hit).
It is re-run from the record and differentiated by reverse-mode AD:
\begin{itemize}
\item hit time: $\phi(x + t v; c) = 0$, so $\partial t^h/\partial c_j = -B_j(x^h)/(v\cdot\nabla\phi(x^h))$
  (bisection for the root, one Newton step that carries the derivative);
\item then $x^h = x + t^h v$, $n(x^h; c)$, $\vout = R_n \vin$,
  $x^\pth = x^h + \vout(\Delta t - t^h)$, repeated for every bounce of the step;
\item grazing hits: $1/(v\cdot\nabla\phi)$ is floored at $\epsilon\,|v|\,|\nabla\phi|$ in the
  derivative only (the primal hit is exact): bounded per-hit weights, bias $O(\epsilon)$.
\end{itemize}
The replay must reproduce the record: any flipped discrete outcome is an error.
\codenote{\code{shape.replay} (per-hit and timing rewards, replay check; \code{REPLAY\_TOL}
$= 10^{-9}$ in \code{ensemble\_mesh}), \code{shape.wall\_map}, \code{wall\_adjoint.advance(trace=True)},
\code{levelset.hit\_body}, \code{levelset.regularised\_quotient} (config \code{grazing} $= \epsilon = 0.05$).}
\end{frame}

% ---------------------------------------------------------------------------------------
\section{Terms}

\begin{frame}{Term 1: transport of the perturbed state (pathwise)}
\[
\text{pathwise}_k = \sum_{i\in\text{surv}}
\nabla_x\lambda^m_k(x^\pth_i, v^\pth_i)\cdot\partial_\theta x^\pth_i
+ \nabla_v\lambda^m_k(x^\pth_i, v^\pth_i)\cdot\partial_\theta v^\pth_i .
\]
\begin{itemize}
\item $\partial_\theta x^\pth$, $\partial_\theta v^\pth$: the VJP of the wall map; nonzero only
  for particles that hit the body in step $k$, so the slopes are fitted at those states only.
\item $\nabla_v\lambda$: reflected velocities.  On a specular body $\partial_\theta\vout$ is
  $O(1)$ per hit; the slope says how the particle's future depends on where it goes.
\item $\nabla_x\lambda$: shifted hit points and hit times.  Its component along the
  particle's own velocity comes from an identity, not from the fit (next frame).
\item Only the survivors of the step contribute
  ($a_k\wedge a^m_k\wedge\neg\,\mathrm{born}_k$).
\end{itemize}
\presonly{\medskip For a convex body in free-molecular flow a reflected particle never
re-hits, so this term is the blockage (flow-field) feedback; with collisions it carries the
gas's response to the deflected particles.}
\codenote{Cotangents from \code{node\_gradients} on \code{(x\_post, v\_post)} in
\code{\_gradient\_step}; \code{hit\_fits=True} (\code{\_body\_hit\_slots}; the same result to
the last bit at $\sim$1/200 of the fitting); rows \code{ShapeTerms.pathwise\_x},
\code{pathwise\_v}.}
\end{frame}

\begin{frame}{$\nabla_x\lambda$ along the flight: the along-flight identity}
Free flight carries $(x - s v, v)$ at time $t$ to $(x, v)$ at $t + s$.  So
\[
\lambda_t(x - s v, v) = \lambda_{t+s}(x, v) + s\,(\Lstar\lambda)(x,v) + O(s^2),
\qquad
\lambda_{t+s}(x,v) = \lambda_t(x,v) + s\,T(x,v) + O(s^2),
\]
$T$ = the change of the objective per unit delay of the particle's whole future (the hits
that cross the window edges), and $s\,\Lstar\lambda$ = the collisions during the extra time.
Hence
\[
v\cdot\nabla_x\lambda = -T - \Lstar\lambda .
\]
\begin{itemize}
\item $T$ is the value of the \emph{timing rewards} (Term 5's band rewards), carried by a
  second value column; $\Lstar\lambda \approx c_{k,i}/\Delta t$ from the collision pull-back.
\item The fitted $\nabla_x\lambda$ keeps its perpendicular part; its component along $v$ is replaced:
\[
\nabla_x\lambda \leftarrow \nabla_x\lambda - \frac{v\cdot\nabla_x\lambda + T + c_{k}/\Delta t}{|v|^2}\,v .
\]
\item Accuracy: NTC gives a particle one collision chance per step wherever it sits in its
  cell, so in the discrete scheme the $\Lstar$ part holds to $O(\nu\Delta t)$; exact without collisions.
\end{itemize}
\codenote{\code{shape.along\_flight(gx, v, t)} with \code{t = lam\_timing + collision / dt}
in \code{\_gradient\_step}; \code{lam\_timing} = the second value column (timing rewards
\code{Replay.timing\_reward}); \code{collision} = \code{Backward.collision}
(\code{gradient\_variants(collision=)}; omitting it in a collisional run drops $\Lstar\lambda$).}
\end{frame}

\begin{frame}{Term 2: the hit reward (tally)}
The value $\lambda^m_k$ starts \emph{after} the move, so the rewards the move itself pays --- the
body hits of step $k$ --- enter the backward by hand, $\mu_k = \lambda^m_k + \sum_j r_{k,j}$; their
derivative is this term, the value's is Term 1.
The reward $r = \frac{F_{\mathrm{num}}}{W\Delta t}\,m(\vin - \vout)_x$ depends on $\theta$
through $\vout = R_n\vin$ (and through $\vin$ of later bounces in the same step).  The
surface tallies are outputs of the wall map, so this is one more cotangent:
\[
\text{hit}_k = \frac{\partial D}{\partial T_k}\cdot\partial_\theta T_k,\qquad
\frac{\partial D}{\partial T_k} = \frac{F_{\mathrm{num}}}{W\Delta t}\,\ind[k\in\text{window}]
\ \text{on the body's } x\text{-momentum}.
\]
Exact given the hits; it needs no fitted field.

\medskip
\textbf{Example}: specular cylinder of radius $r$, free-molecular.  With impact parameter
$b$, $F \propto \int_{-r}^{r}\cos^2\varphi\,\dd b = 4r/3$; differentiating at fixed rays gives
$\int 2b^2/r^3\,\dd b = 4/3$, so this term alone is $F/r$ --- the whole derivative of an
isolated body, since the new grazing rays carry $\cos^2\varphi = 0$.  The other terms carry
the response of the flow (blockage, collisions).
\codenote{\code{ensemble\_mesh.\_drag\_tally\_cotangent} $\times$ \code{tally\_weight} in
\code{\_gradient\_step}; row \code{ShapeTerms.pathwise\_tally}.}
\end{frame}

\begin{frame}{Term 3: diffuse walls, the normal of the re-emission law}
A diffuse hit re-emits from $q_w(v\mid n) = \frac{v_n}{s^2}e^{-v_n^2/2s^2}\,\frac{e^{-|v_t-u|^2/2s^2}}{2\pi s^2}$,
$s^2 = kT_w/m$.  Its normal moves with $\theta$.  Split the hit's future into its reward and
the value after it:
\[
\frac{\dd}{\dd n}\E\big[r(\vout) + V(\vout)\big]
= \underbrace{\E\Big[\frac{\partial r}{\partial \vout}\cdot\frac{\partial \vout}{\partial n}\Big]}_{\text{pathwise, reparameterised}}
+ \underbrace{\E\big[V(\vout)\,\partial_n\log q_w(\vout\mid n)\big]}_{\text{score, value after the hit}} .
\]
\begin{itemize}
\item Pathwise part: $\vout = s\,(\rho\, n + z_1 t_1(n) + z_2 t_2(n)) + u$ at fixed draws;
  the tangent frame is parallel-transported from the primal one (bounded derivative), and
  the reward enters through the tally (Term 2).
\item Score part: weighted by the value-to-go \emph{after} the hit; no fitted slope.
  $1/v_n$ is floored at $\epsilon|\vout|$ (bias $\approx\epsilon/2$ of the $v_t/v_n$ part).
\item Specular bodies: zero.
\end{itemize}
\codenote{\code{wall\_adjoint.\_element\_events} (tally from the traced sample),
\code{sampling.diffuse\_velocity\_general} (parallel-transported frame),
\code{shape.normal\_score}, \code{sampling.log\_diffuse\_density(floor=)}; weights
\code{to\_go.post} when a tally cotangent is given; row \code{ShapeTerms.score}.}
\end{frame}

\begin{frame}{Term 4: hit $\leftrightarrow$ miss (silhouette band)}
Along a particle's undeflected line, $\phi_{\min} = \min_s\phi(x + s\hat v)$: $<0$ it hits,
$>0$ it misses.  A surface displacement flips the outcome where $|\phi_{\min}| < |\delta\phi|$:
\[
\frac{\dd D}{\dd\theta}\Big|_{\mathrm{sil}} = -\E\Big[\rho(0)\,(V_{\mathrm{hit}} - V_{\mathrm{miss}})\,\frac{\partial\phi_{\min}}{\partial\theta}\Big]
\approx -\frac{1}{2\varepsilon_b}\sum_{|\phi_{\min,i}|<\varepsilon_b}(V_{\mathrm{hit},i} - V_{\mathrm{miss},i})\,B_j(x^*_i)\,\frac{\partial c_j}{\partial\theta},
\]
$\partial\phi_{\min}/\partial c_j = B_j(x^*)$ at the closest point (envelope theorem).
\begin{itemize}
\item Realised outcome: its reward-to-go.  Virtual outcome: the fitted value at the virtual
  post-move state plus the virtual hit's reward, valued from the leaving population.
\item Each passing counted once per trajectory (a band wider than one move would otherwise
  count it several times).
\item Specular: $V_{\mathrm{hit}}\to V_{\mathrm{miss}}$ at grazing incidence, the term is $O(\varepsilon_b)$
  and is left out.  Diffuse: $O(1)$ --- about half of a diffuse body's derivative.
\end{itemize}
\codenote{\code{shape.silhouette\_cotangent}, \code{levelset.closest\_approach},
\code{ensemble\_device.value\_query}; config \code{silhouette\_band} $=\varepsilon_b$ (0 for
specular bodies, 0.03 on the diffuse squash); row \code{ShapeTerms.silhouette}.}
\end{frame}

\begin{frame}{Term 5: hits that cross a window edge}
The VJP differentiates a fixed assignment of hits to steps.  A surface displacement shifts
hit times, and a hit at the end of step $k$ can slip into step $k+1$.  For the per-hit drag
this changes the objective only at the window's edges:
\[
\Delta J_{k,j} = +r_{k,j}\ \text{(last window step)},\qquad -r_{k,j}\ \text{(step before the window)} .
\]
Boundary term: (density of hit times at the step end) $\times(-\Delta J)$, the density estimated
over the hits of the $w$ whole steps ending at each edge:
\[
\text{cotangent on } t^h_{k,j}:\qquad -\frac{\pm r_{k,j}/w}{\Delta t}
\quad\text{for every body hit of those steps.}
\]
\begin{itemize}
\item Zero in expectation on a steady window, but part of the windowed objective's derivative;
  it is included.  $w = 10$ steps: $\sim$50$\times$ the hits of a fraction-of-a-step band, at an
  $O(w\Delta t)$ smoothing of the edge's hit-time density.
\item The same per-hit band rewards are the \emph{timing rewards} whose value $T$ enters the
  along-flight identity.
\end{itemize}
\codenote{\code{ensemble\_mesh.\_drag\_jump(width=w)}, \code{timing\_band\_setting} (a value
$\ge 1$ = whole steps per edge; default \code{timing\_band} $= 10$),
\code{shape.timing\_cotangent}; row \code{ShapeTerms.timing}.}
\end{frame}

\begin{frame}{Term 6: the order of a hit and the step's collisions}
NTC collides once per step, after the move.  A body hit at the end of step $k$ that slips into
step $k+1$ is, during step $k$'s collisions, the \emph{incoming} particle $(x_w, \vin)$
instead of the \emph{reflected} one $(x_w, \vout)$.  Both futures coincide after the next
move, so the objective changes by
\[
J_{c} = c(x_w, \vin) - c(x_w, \vout),\qquad c = \Delta t\,\Lstar\lambda ,
\]
at \emph{every} step boundary, not only at the window's edges.  The hit times of a steady flow
are uniform within the step, so the density at the step end is $1/\Delta t$ and
\[
\text{cotangent on } t^h:\qquad +\frac{J_c}{\Delta t}\quad\text{for every body hit.}
\]
\begin{itemize}
\item $c$ at the hit point: fitted from the collision increments $c_{k,i}$ of the hit's
  cell, on the approaching side ($v\cdot n<0$) for $\vin$ and the leaving side for $\vout$.
\item Without collisions $c\equiv0$; then Term 5 is the whole hit-step dependence.
\end{itemize}
\codenote{\code{shape.order\_jump} (\code{ORDER\_FIT} $= \{k: 32,\ \text{pool}: 0\}$;
\code{value\_query(direction=}$\mp n$\code{)} on \code{Backward.collision}); cotangent
\code{timing\_cotangent(head, -jump, dt, 1.0)} on the first pass's body hits; row
\code{ShapeTerms.order} (set when \code{collision} is given).}
\end{frame}

\begin{frame}{Term 7: the cut-cell volume in the collision step}
Every collision probability of cell $c$ scales as $1/V_c$ (NTC frame).  The collision law is
the only thing that changes, so this is a kernel-parameter (operator) term with
$\partial_\theta\log q$ replaced by $\partial_\theta\log V_c^{-1}$:
\[
\frac{\dd D}{\dd\theta}\Big|_V = -\sum_k\sum_c\frac{\partial_\theta V_c}{V_c}\,C_{k,c},
\qquad
C_{k,c} = \tfrac12\sum_{i\in c} c_{k,i}\quad\text{(each pair seen from both ends)},
\]
$C_{k,c}$ the expected change of the objective by the collisions of cell $c$ in step $k$.
The geometry factor is the shape derivative of the cut volume:
\[
V_c = \int_c H(\phi)\,\dd x,\qquad
\frac{\partial V_c}{\partial c_j} = \int_c\delta(\phi)\,B_j\,\dd x
= \int_{\Gamma\cap c}\frac{B_j}{|\nabla\phi|}\,\dd s ,
\]
evaluated with a smoothed Heaviside $H_w(\phi) = \tfrac12\big(1 + \tanh(\phi/w)\big)$; each
sub-region of a split cell separately.  Zero without collisions.
\codenote{\code{Backward.cell\_collision} (per-cell sums of $c_{k,i}$),
\code{shape.volume\_weights} ($g_c = -\sum_k C_{k,c}/V_c$), \code{shape.volume\_cotangent}
($32^D$ samples per cell, $w = h_\phi/4$, split cells by \code{collision\_cell}),
\code{ensemble\_mesh.\_volume\_term} (level-set bodies, virtual pull-back); row \code{volume}.}
\end{frame}

% ---------------------------------------------------------------------------------------
\section{Assembly}

\begin{frame}{Assembly}
Per seed, on the level-set coefficients:
\[
\widehat{\nabla_c D} = \sum_k\Big(\underbrace{\text{pathwise}_x + \text{pathwise}_v}_{1}
+ \underbrace{\text{tally}}_{2} + \underbrace{\text{score}}_{3} + \underbrace{\text{silhouette}}_{4}
+ \underbrace{\text{timing}}_{5} + \underbrace{\text{order}}_{6}\Big) + \underbrace{\text{volume}}_{7},
\]
then chained to the controls by the VJP of $\theta\mapsto c(\theta)$, and averaged over seeds.

\begin{center}\footnotesize
\begin{tabular}{@{}lccc@{}}
\toprule
term & free-molecular, specular & free-molecular, diffuse & collisional\\
\midrule
1 pathwise $x$, $v$ & \ok & \ok & \ok{} ($+\Lstar\lambda$ along $v$)\\
2 tally & \ok & \ok & \ok\\
3 diffuse normal score & \nil & \ok & if diffuse\\
4 silhouette & $O(\varepsilon_b)$, off & \ok & if diffuse\\
5 window edges & \ok & \ok & \ok\\
6 hit/collision order & \nil & \nil & \ok\\
7 cut-cell volume & \nil & \nil & \ok\\
\bottomrule
\end{tabular}
\end{center}

One gradient evaluation: forward with capture $\to$ replay $\to$ two-column value backward
(with the collision increments) $\to$ one VJP scan over the captured steps $\to$ volume term
$\to$ control chain.
\codenote{\code{ensemble\_mesh.shape\_terms} $\to$ \code{ensemble\_gradient(variant=}
\code{"along\_pool\_plane")}: \code{gradient = pathwise + score + silhouette + volume + timing
+ order}; \code{control\_gradient} (\code{jax.vjp} of \code{levelset\_coef}).  A
\code{newborn\_hit\_share} above \code{NEWBORN\_HIT\_MAX} $=10^{-3}$ warns (\code{\_newborn\_share}).}
\end{frame}

\begin{frame}{From the gradient to the optimiser}
\textbf{Squash (one scalar control).}  Adam on
$J'(s) = \widehat{\dd D/\dd s} - \lambda_A\,\dd A/\dd s$, the step's largest change of $s$ set by
the learning rate; $S$ common-random-number seeds per iterate.

\medskip
\textbf{Free shape (all coefficients), at fixed area.}
\begin{enumerate}
\item seed mean $g$; restrict to the coefficients within a band of the surface and smooth with a
  Gaussian (an $H^1$-type gradient that removes knot-scale noise);
\item gradient descent or Adam (Adam's normalised step is smoothed again);
\item fixed largest normal displacement per iterate ($s_{\max} = 0.05$, a tenth of a cell);
\item reinitialise $\phi$ to a signed distance;
\item restore the area by a uniform shift $\gamma$ of $\phi$: to first order the displacement becomes
  $V_n = \eta\,(d - \langle d\rangle_\Gamma)$, the projection of the step onto
  $\{\oint V_n = 0\}$, and the constraint then holds exactly (a retraction).
\end{enumerate}
\codenote{\code{optimize\_squash.run}, \code{seed\_gradient}, \code{adam\_direction}
(\code{learning\_rate} $= 0.05$); \code{optimize.run}, \code{smoothed\_direction}
($\sigma = 2$ knots, band $1 + h_\phi$), \code{levelset.reinitialise},
\code{project\_area} (bisection on \code{levelset.body\_volume}).  Both default to
\code{along\_pool\_plane} and \code{timing\_band} $= 10$.}
\end{frame}

\begin{frame}{What is exact, and what is approximate}
\small
\begin{tabular}{@{}p{4.0cm}p{5.2cm}p{3.9cm}@{}}
\toprule
piece & approximation & size / control\\
\midrule
wall-map VJP, tally, scores & exact given the record & ---\\
collision pull-back & $S$ substeps, $K$ virtual partners & $O(\Delta t/S)$; $(S,K)=(4,8)$\\
value field of the pull-back & one guarded quadratic per cell (split at surfaces) & $-3.7\pm0.8\,\%$ (Kn 0.5)\\
fitted slopes $\nabla\lambda$ & local quadratic, state-based node sets & fit bias; $k$, pool, half-space\\
along-flight $\Lstar$ part & continuous identity in a discrete scheme & $O(\nu\Delta t)$\\
grazing floor & $1/(v\cdot n)$ floored in the derivative & $O(\epsilon)$; $\epsilon = 0.05$\\
silhouette band & finite band density & $O(\varepsilon_b)$\\
window-edge band & $w$ whole steps & $O(w\Delta t)$ smoothing\\
order term & uniform hit times in the step & steady flow\\
captured horizon & sensitivities older than the captured steps & box memory (walls)\\
\bottomrule
\end{tabular}

\medskip
\textbf{Left out on purpose}: hits by particles born in the step or leaving the box in it
(monitored, $0$ while the body is away from the inflow and outflow faces); the change of cell
membership of a hit particle ($O(\nu\Delta t)$ below the order term); the silhouette term of
specular bodies.

\medskip
Validation of each piece against finite differences: the validation deck.
\codenote{Settings: \code{ensemble\_backward.VIRTUAL}; \code{shape.VARIANTS};
\code{grazing}; \code{silhouette\_band}; \code{timing\_band}; captured \code{steps}
(300 in production, $\ge 600$ for 1\,\%-level free-molecular gates); \code{NEWBORN\_HIT\_MAX}.}
\end{frame}
